% kap03.tex — Hamilton-formalisme
% Hamilton-Jacobi-kompendiet

\chapter{Hamilton-formalisme}
\label{kap:hamilton}

\section{Målet: fra hastigheder til impulser}

I kapitel~\ref{kap:lagrange} stødte vi to gange på den generaliserede
impuls $p_i=\partial L/\partial\dot q_i$ (afsnit om cykliske
koordinater) og på størrelsen
$H=\sum_i\dot q_i\,\partial L/\partial\dot q_i - L$
(afsnit om Noethers sætning), uden at give dem en systematisk
behandling. I dette kapitel gør vi netop det: vi bytter de $n$
generaliserede hastigheder $\dot q_i$ ud med de $n$ generaliserede
impulser $p_i$ som de fundamentale variable, og omskriver hele
dynamikken i termer af $(q_i,p_i)$ i stedet for $(q_i,\dot q_i)$.

Det matematiske redskab til denne ombytning er
\emph{Legendre-transformationen}\index{Legendre-transformation}. Resultatet er \emph{Hamiltons
ligninger}\index{Hamiltons ligninger} — to \emph{første}ordens differentialligninger for hver
frihedsgrad, i stedet for Euler-Lagranges \emph{ene anden}ordens
ligning. Denne omskrivning er ikke blot kosmetisk: den flytter
beskrivelsen fra konfigurationsrummet ($n$-dimensionalt, koordinaterne
$q_i$) til \emph{faserummet}\index{faserum} ($2n$-dimensionalt, $(q_i,p_i)$), som er
selve arenaen for al den resterende teori i kompendiet — kanoniske
transformationer, Poisson-parenteser og til sidst Hamilton-Jacobi-
ligningen selv.

\section{Legendre-transformationen}

\subsection{Den generelle idé, illustreret i én variabel}

Før vi anvender Legendre-transformationen på $L$, er det nyttigt at se
den i sin enkleste form. Lad $f(x)$ være en konveks funktion af én
variabel $x$ (dvs.\ $f''(x)>0$ overalt). Vi definerer en ny variabel
\begin{equation}
  p \equiv \dv{f}{x},
\end{equation}
og — fordi $f$ er konveks, er $p(x)$ monotont voksende og dermed
invertibel: vi kan i princippet løse for $x$ som funktion af $p$. Den
\emph{Legendre-transformerede} funktion defineres som
\begin{equation}
  g(p) \equiv x\,p - f(x), \qquad \text{hvor } x = x(p) \text{ (den inverse funktion).}
  \label{eq:legendre-generel}
\end{equation}
Pointen er, at $g$ indeholder \emph{præcis} den samme information som
$f$ — vi kan altid gå tilbage — men udtrykt i den nye variabel $p$ i
stedet for $x$. At vi ganger med $p$ og trækker $f$ fra, er ikke en
arbitrær opskrift: det er netop den kombination, der garanterer, at
$dg/dp = x$ (vis selv dette ved at differentiere
\eqref{eq:legendre-generel} med hensyn til $p$ og bruge
$dp/dx\cdot dx/dp=1$), så transformationen er sin egen invers, op til
et fortegn — en symmetri, vi skal udnytte om et øjeblik.

\begin{bemaerkning}
Legendre-transformationen er langt fra unik for mekanik. Den samme
konstruktion optræder i termodynamikken (fra intern energi $U(S,V)$
til fri energi $F(T,V)=U-TS$, hvor $T=\partial U/\partial S$) og i
konveks optimering. Den fælles struktur er altid: erstat en variabel
med dens ``partielle afledte-partner'', og betal for skiftet med et
produktled.
\end{bemaerkning}

\subsection{Anvendt på Lagrange-funktionen}

Vi anvender nu \eqref{eq:legendre-generel} led for led på
$L(q_1,\dots,q_n,\dot q_1,\dots,\dot q_n,t)$, med de $n$ hastigheder
$\dot q_i$ som de variable, der skal udskiftes. For hver $i$ definerer
vi den \emph{kanoniske impuls}\index{kanonisk impuls}
\begin{equation}
  \boxed{p_i \equiv \pdv{L}{\dot q_i}.}
  \label{eq:def-p}
\end{equation}
For at denne transformation er veldefineret — dvs.\ at vi entydigt kan
løse \eqref{eq:def-p} for $\dot q_i$ som funktion af $(q,p,t)$ — skal
Hesse-matricen
\begin{equation}
  \left(\pdv[2]{L}{\dot q_i}{\dot q_j}\right)
\end{equation}
være ikke-singulær (invertibel implicit funktionssætning). Dette er
opfyldt for stort set alle fysisk relevante Lagrange-funktioner —
konkret er det opfyldt, når den kinetiske energi $T$ er en ægte
kvadratisk (positivt definit) funktion af hastighederne, hvad den er
for alle systemer, vi møder i dette kompendium.

\begin{bemaerkning}
Systemer, hvor Hesse-matricen \emph{er} singulær, kaldes
``tvangssystemer'' (constrained systems) i en teknisk betydning, der er
snævrere end de holonome tvangsbetingelser fra
kapitel~\ref{kap:newton-variation} — de kræver Diracs teori for
tvangssystemer. Dette optræder typisk for felter med gaugesymmetri
(elektromagnetisme, Yang-Mills-teori), og bliver derfor relevant, når
vi i kapitel~\ref{kap:felter} og i det senere QFT-kompendium når til
gaugefelter. Vi nævner det her, så du genkender det, når det opstår —
men det ligger uden for dette kompendiums rækkevidde at behandle det.
\end{bemaerkning}

Vi definerer nu \emph{Hamilton-funktionen}\index{Hamilton-funktion} som den
Legendre-transformerede af $L$ med hensyn til alle $\dot q_i$
samtidigt:
\begin{equation}
  \boxed{H(q,p,t) \equiv \sum_{i=1}^n p_i\dot q_i - L(q,\dot q,t),}
  \label{eq:def-H}
\end{equation}
hvor $\dot q_i$ på højresiden skal udtrykkes ved $(q,p,t)$ via den
(lokale) inversion af \eqref{eq:def-p}. Dette er netop den størrelse,
vi kaldte $H$ i \eqref{eq:def-H-forlobig} i kapitel~\ref{kap:lagrange}
— nu givet en fuldstændig, systematisk definition som funktion af
$(q,p,t)$ i stedet for $(q,\dot q,t)$.

\section{Hamiltons ligninger}

\subsection{Udledning ved direkte differentiation}

Vi differentierer definitionen \eqref{eq:def-H} totalt:
\begin{equation}
  dH = \sum_i\left(\dot q_i\,dp_i + p_i\,d\dot q_i\right) - dL.
\end{equation}
Da $L=L(q,\dot q,t)$, er
\begin{equation}
  dL = \sum_i\left(\pdv{L}{q_i}dq_i + \pdv{L}{\dot q_i}d\dot q_i\right) + \pdv{L}{t}dt
     = \sum_i\left(\pdv{L}{q_i}dq_i + p_i\,d\dot q_i\right) + \pdv{L}{t}dt,
\end{equation}
hvor vi brugte definitionen \eqref{eq:def-p} i sidste skridt. Indsat i
udtrykket for $dH$ ovenfor ser vi, at leddene med $d\dot q_i$
\emph{går ud} — det er netop pointen ved en Legendre-transformation:
\begin{equation}
  dH = \sum_i\left(\dot q_i\,dp_i - \pdv{L}{q_i}dq_i\right) - \pdv{L}{t}dt.
  \label{eq:dH-mellemresultat}
\end{equation}
Vi genkender nu $\partial L/\partial q_i$ fra Euler-Lagrange-ligningen
\eqref{eq:euler-lagrange}: langs en faktisk løsning er
$\partial L/\partial q_i = \dv*{}{t}(\partial L/\partial \dot q_i) =
\dot p_i$. Med denne substitution bliver
\eqref{eq:dH-mellemresultat}:
\begin{equation}
  dH = \sum_i\left(\dot q_i\,dp_i - \dot p_i\,dq_i\right) - \pdv{L}{t}dt.
  \label{eq:dH-slutresultat}
\end{equation}

På den anden side er $H=H(q,p,t)$ per konstruktion, så den totale
differential \emph{også} kan skrives
\begin{equation}
  dH = \sum_i\left(\pdv{H}{q_i}dq_i + \pdv{H}{p_i}dp_i\right) + \pdv{H}{t}dt.
  \label{eq:dH-fra-H}
\end{equation}
Da $q_i$, $p_i$ og $t$ er uafhængige variable i $H$'s beskrivelse, må
koefficienterne til $dq_i$, $dp_i$ og $dt$ i
\eqref{eq:dH-slutresultat} og \eqref{eq:dH-fra-H} stemme overens led
for led. Dette giver os \emph{Hamiltons ligninger}:
\begin{equation}
  \boxed{\dot q_i = \pdv{H}{p_i}, \qquad \dot p_i = -\pdv{H}{q_i}, \qquad i=1,\dots,n,}
  \label{eq:hamiltons-ligninger}
\end{equation}
samt, som en biprodukt af samme sammenligning,
\begin{equation}
  \pdv{H}{t} = -\pdv{L}{t}.
  \label{eq:dH-dt}
\end{equation}

\begin{bemaerkning}
Hamiltons ligninger \eqref{eq:hamiltons-ligninger} er $2n$ \emph{først}e-
ordens differentialligninger, mod Euler-Lagranges $n$ \emph{anden}-
ordens ligninger. Dette er ikke en tilfældig gevinst: en tilstand i
faserummet, $(q_i(t),p_i(t))$, er nu fuldstændig bestemt af sin værdi
ved \emph{ét} tidspunkt — der er ikke brug for at kende både position
\emph{og} hastighed ved starttidspunktet særskilt, fordi $p_i$ allerede
``indeholder'' den information, der før lå i $\dot q_i$. Denne
førsteordens-struktur er præcis, hvad der gør faserumsbilledet så
naturligt, og den er en direkte forløber for tilstandsbegrebet i
kvantemekanikken, hvor en tilstand (bølgefunktionen) også er
fuldstændigt bestemt ved et enkelt tidspunkt.
\end{bemaerkning}

\subsection{Fortolkningen af $H$: bevarelse og energi}

Vi genbesøger nu \eqref{eq:dH-dt} i lyset af Hamiltons ligninger. Den
totale tidsafledte af $H$ langs en faktisk bane er
\begin{equation}
  \dv{H}{t} = \sum_i\left(\pdv{H}{q_i}\dot q_i + \pdv{H}{p_i}\dot p_i\right) + \pdv{H}{t}
  = \sum_i\left(\pdv{H}{q_i}\pdv{H}{p_i} - \pdv{H}{p_i}\pdv{H}{q_i}\right) + \pdv{H}{t}
  = \pdv{H}{t},
\end{equation}
hvor vi brugte Hamiltons ligninger \eqref{eq:hamiltons-ligninger} til
at indsætte $\dot q_i$ og $\dot p_i$, og de to led i parentesen
ophæver hinanden identisk. Vi har altså det generelle resultat
\begin{equation}
  \dv{H}{t} = \pdv{H}{t}.
  \label{eq:dHdt-generel}
\end{equation}
Specielt: hvis $H$ ikke afhænger \emph{eksplicit} af tiden
($\partial H/\partial t=0$, hvilket ifølge \eqref{eq:dH-dt} er
ækvivalent med $\partial L/\partial t=0$), er $H$ en \emph{bevaret
størrelse} langs bevægelsen — vi genfinder Noether-resultatet fra
kapitel~\ref{kap:lagrange}, nu udledt direkte i Hamilton-billedet, uden
omvejen over $L$.

For de systemer, vi typisk møder — tidsuafhængige holonome
tvangsbetingelser, og et potential $V$, der ikke afhænger af
hastigheder — viste opgave 3 i kapitel~\ref{kap:lagrange}, at
$H=T+V$, den samlede mekaniske energi. I disse tilfælde er Hamiltons
ligninger derfor blot en omskrivning af energibevarelse og
bevægelsesligningerne i førsteordensform — men formalismens virkelige
styrke viser sig, når $H\neq T+V$ (fx i tidsafhængige eller
elektromagnetiske systemer), hvor $H$ stadig er den korrekte
``generator'' af tidsudviklingen, selv om den ikke er den mekaniske
energi.

\section{Eksempel: den harmoniske oscillator i faserum}

\begin{eksempel}
For en énddimensional harmonisk oscillator med masse $m$ og
fjederkonstant $k$ er
\begin{equation}
  L = \frac12 m\dot x^2 - \frac12 kx^2.
\end{equation}
Den kanoniske impuls er
\begin{equation}
  p = \pdv{L}{\dot x} = m\dot x
  \quad\Longrightarrow\quad
  \dot x = \frac{p}{m},
\end{equation}
som her simpelthen er den velkendte bevægelsesmængde — Legendre-
transformationen er triviel at invertere for dette system. Hamilton-
funktionen bliver
\begin{equation}
  H = p\dot x - L = p\cdot\frac{p}{m} - \left(\frac12 m\left(\frac{p}{m}\right)^2 - \frac12 kx^2\right)
    = \frac{p^2}{2m} + \frac12 kx^2,
\end{equation}
altså netop $T+V$ udtrykt ved $p$ i stedet for $\dot x$, som forventet.
Hamiltons ligninger \eqref{eq:hamiltons-ligninger} giver
\begin{equation}
  \dot x = \pdv{H}{p} = \frac{p}{m}, \qquad
  \dot p = -\pdv{H}{x} = -kx.
\end{equation}
Den første ligning er blot definitionen af $p$ igen (som det altid vil
være for et system uden hastighedsafhængigt potential); den anden er
Newtons anden lov, $m\ddot x=-kx$, skrevet i førsteordensform. Løsningen
er velkendt, $x(t)=A\cos(\Omega t+\varphi)$ med $\Omega=\sqrt{k/m}$, og
banen i faserummet $(x,p)$ er ellipsen
\begin{equation}
  \frac{p^2}{2mH} + \frac{x^2}{2H/k} = 1,
\end{equation}
fordi $H$ er konstant langs banen (se figur~\ref{fig:faseportraet}).
Dette er vores første eksempel på et \emph{faseportræt}\index{faseportræt}: en geometrisk
repræsentation af systemets dynamik som en kurve (eller familie af
kurver, for forskellige energier $H$) i faserummet, uden nogen
eksplicit reference til tiden $t$.
\end{eksempel}

\begin{figure}[htbp]
  \centering
  \begin{tikzpicture}[scale=1.0]
    % Akser
    \draw[-{Latex}] (-3.2,0) -- (3.2,0) node[right] {$x$};
    \draw[-{Latex}] (0,-2.4) -- (0,2.4) node[above] {$p$};
    % Familie af faseportræt-ellipser for forskellige energier
    \foreach \sc/\col in {0.5/gray!70,0.8/gray!45,1.1/black} {
      \draw[thick,\col] (0,0) ellipse ({2.6*\sc} and {1.9*\sc});
    }
    % Pil der viser bevægelsesretning (med uret) på yderste ellipse
    \draw[-{Latex},thick] ({2.6*1.1*cos(20)},{1.9*1.1*sin(20)}) -- ({2.6*1.1*cos(-5)},{1.9*1.1*sin(-5)});
    \node at (2.9,1.6) {$H_3$};
    \node at (2.15,1.2) {$H_2$};
    \node at (1.4,0.75) {$H_1$};
  \end{tikzpicture}
  \caption{Faseportræt for den harmoniske oscillator: hver ellipse er en
  bane med konstant energi $H$ (voksende fra $H_1$ til $H_3$).
  Bevægelsen langs hver ellipse sker med uret, som markeret med pilen.
  Faserumsbeskrivelsen erstatter tidsforløbet $x(t)$ med en statisk,
  geometrisk figur.}
  \label{fig:faseportraet}
\end{figure}

\section{Faserummet som arena}

Det er værd at fremhæve, hvad vi har vundet ved skiftet til
$(q_i,p_i)$. I konfigurationsrummet ($q_i$ alene) skærer baner for
forskellige startbetingelser hinanden hele tiden — en partikel, der
passerer gennem et bestemt punkt $x_0$, kan have en hvilken som helst
hastighed, så uendeligt mange baner går gennem samme punkt. I
\emph{faserummet} $(q_i,p_i)$ er dette udelukket: gennem hvert punkt
$(q_0,p_0)$ går \'en og kun \'en bane, fordi $(q_0,p_0)$ er en
fuldstændig starttilstand. Faserumsbaner for forskellige startbetingelser
kan derfor aldrig skære hinanden (bortset fra i entydige ligevægtspunkter)
— en egenskab, der bliver central, når vi i kapitel~\ref{kap:kanoniske-transf}
og \ref{kap:hj} studerer, hvordan hele familier af baner transformerer
sig samtidigt under kanoniske transformationer.

\begin{bemaerkning}
Faserummets rolle som en arena, hvori ``tilstanden'' er et punkt, og hvor
baner ikke skærer hinanden, er selve det billede, der (i en passende
kvantiseret og ``udtværet'' form, jf.\ Heisenbergs
usikkerhedsrelation) bliver til kvantemekanikkens tilstandsrum. Vejen
fra det klassiske faserum til kvantemekanikkens Hilbert-rum går netop
via de strukturer, vi udvikler i de kommende kapitler — Poisson-
parenteser (kapitel~\ref{kap:poisson}), som bliver til kommutatorer ved
kanonisk kvantisering, og Hamilton-Jacobi-ligningen
(kapitel~\ref{kap:hj}), som viser sig at være den klassiske grænse af
Schrödinger-ligningen (se [[kvantemekanik-kompendiet]] for den
kvantemekaniske side af denne forbindelse).
\end{bemaerkning}

\section{Opsummering}

\begin{itemize}
  \item Legendre-transformationen udskifter en variabel med dens
        partielle afledte-partner: $g(p)=xp-f(x)$, hvor $p=df/dx$.
  \item Anvendt på $L(q,\dot q,t)$ med $p_i=\partial L/\partial\dot q_i$
        giver dette Hamilton-funktionen
        $H(q,p,t)=\sum_i p_i\dot q_i - L$.
  \item Sammenligning af $dH$ udledt af definitionen med $dH$ udledt af
        $H(q,p,t)$ giver Hamiltons ligninger,
        $\dot q_i=\partial H/\partial p_i$,
        $\dot p_i=-\partial H/\partial q_i$ — $2n$
        førsteordensligninger i stedet for $n$ andenordensligninger.
  \item $dH/dt = \partial H/\partial t$: $H$ er bevaret, når den ikke
        afhænger eksplicit af tiden. For tidsuafhængige,
        hastighedsuafhængige potentialer er $H=T+V$.
  \item Faserummet $(q_i,p_i)$ er arenaen for resten af kompendiet:
        hvert punkt er en fuldstændig tilstand, og baner for forskellige
        starttilstande skærer aldrig hinanden.
\end{itemize}

\begin{opgave}
Opskriv Lagrange-funktionen for en partikel med masse $m$ i et
centralt potential $V(r)$ i planpolar-koordinater $(r,\phi)$ (jf.\
opgave 2 i kapitel~\ref{kap:lagrange}), udfør Legendre-transformationen
til $(p_r,p_\phi)$, og vis at
\begin{equation}
  H = \frac{p_r^2}{2m} + \frac{p_\phi^2}{2mr^2} + V(r).
\end{equation}
Opskriv Hamiltons ligninger for $(r,p_r,\phi,p_\phi)$ og sammenlign med
resultatet fra kapitel~\ref{kap:lagrange}.
\end{opgave}

\begin{opgave}
En ladet partikel med masse $m$ og ladning $e$ bevæger sig i et
elektromagnetisk felt beskrevet ved skalarpotentialet $\Phi$ og
vektorpotentialet $\vb{A}$. Lagrange-funktionen er (uden bevis her —
den udledes korrekt først i forbindelse med feltteorien i
kapitel~\ref{kap:felter})
\begin{equation}
  L = \frac12 m\dot{\vb{r}}^2 - e\Phi + e\,\dot{\vb{r}}\cdot\vb{A}.
\end{equation}
Vis, at den kanoniske impuls er $\vb{p}=m\dot{\vb{r}}+e\vb{A}$ — altså
\emph{ikke} lig med den mekaniske bevægelsesmængde $m\dot{\vb{r}}$ —
og find Hamilton-funktionen $H(\vb{r},\vb{p},t)$. Dette er det
kanoniske eksempel på, at $H$ og den mekaniske energi $T+V$ kan være to
forskellige størrelser.
\end{opgave}

\begin{opgave}
Skitser (i hånden eller med et plot) faseportrættet for en partikel med
energi $E$, der bevæger sig i et énddimensionalt potential med en
enkelt potentialbrønd (fx $V(x)=\frac12 kx^2$ for $|x|<a$ og
$V(x)=\infty$ for $|x|\geq a$, en ``boks med fjederbund''). Angiv, hvilke
værdier af $E$ der giver lukkede (periodiske) baner, og diskuter, hvad
der sker med faseportrættet, når $E$ nærmer sig den energi, hvor
partiklen netop når randen $|x|=a$.
\end{opgave}
